%% main.tex -- the letter of BRIEF_v2.md (17 Nov 2026).
%% Every number is a macro: values/results_macros.tex (shared), values_independent_results.tex
%% (real results, generated) and values_independent.tex (expected values, red). A sentence whose
%% wording depends on how a run comes out ends with the comment EXPECTED.
\documentclass[aps,prd,reprint,amsmath,amssymb,nofootinbib]{revtex4-2}
\usepackage{graphicx}
\usepackage{xcolor}
\usepackage{hyperref}
\hypersetup{colorlinks=true,linkcolor=[rgb]{0.10,0.20,0.55},citecolor=[rgb]{0.10,0.20,0.55},urlcolor=[rgb]{0.10,0.20,0.55}}

\input{values/results_macros.tex}
\InputIfFileExists{values_independent_results.tex}{}{}
%% values_independent.tex -- numbers of the short paper that values/results_macros.tex does not carry.
%% Input AFTER values/results_macros.tex. Every macro is \providecommand, so a value file that
%% defines the same name earlier wins. Red (\ph): expected value, to be replaced when the
%% experiment of EXPERIMENT.md finishes. Black: design value of the simulation.
%% Meaning, expected value and its source for every macro: notes/MACROS_NEEDED.md.
%% Half-widths are half the width of the 90% credible interval of H0, km/s/Mpc, 1000 mergers.
%% 17 Nov 2026: where stage 1 of the production list has measured the same kind of analysis and the macro's own
%% run of the experiment has not finished, the red value is the stage-1 measurement (marked MEASURED below).

%% ---- design values not in values/MACROS.md (black) ----
\providecommand{\IndNreal}{five}         % realizations of the experiment (a word: use in text mode)
\providecommand{\IndfList}{0,\ 0.1,\ 0.3,\ 0.5\ {\rm and}\ 1}  % the input fractions of Fig. 2 (math mode)
\providecommand{\IndDmTwice}{10}         % mass-peak difference of the stronger case [Msun]
\providecommand{\IndDchiTwice}{0.20}     % mean effective spin difference of the stronger case

%% ---- Fig. 1 (figures/fig_concept, drawn from the simulated catalogs; final, black) ----
\providecommand{\IndConceptDec}{3}             % half-thickness of the slice in declination [degrees]
\providecommand{\IndConceptzMax}{0.20}         % largest redshift of the slice
\providecommand{\IndConceptScale}{200}         % length of the scale bar [Mpc, comoving]
\providecommand{\IndConceptNGal}{13{,}398}     % objects of the galaxy catalog in the slice
\providecommand{\IndConceptNAGN}{126}          % objects of the AGN catalog in the slice
\providecommand{\IndConceptzTrue}{0.098}       % true redshift of the merger whose 90% region is outlined
\providecommand{\IndConceptVolNGal}{225}       % galaxies inside that 90% localization region
\providecommand{\IndConceptVolNAGN}{6}         % AGN inside that 90% localization region
\providecommand{\IndConceptNgalHosted}{640}    % galaxy-hosted among the 1000 detected mergers, populations differ
\providecommand{\IndConceptNAGNHosted}{360}    % AGN-hosted among the 1000 detected mergers, populations differ

%% ---- Fig. 2(a) (figures/fig1_budget): sharp mass scale, populations held at their inputs (red) ----
\providecommand{\IndHWSpec}{\ph{6.0}}          % mass spectrum alone, no catalog. MEASURED in stage 1 on the f = 0.30 sets (known_population/no_catalog/complete/equal_populations): mean over realizations 100, 101, 102, 103, 105 is 6.03 (range 5.95 to 6.14; realization A 6.06). Was 3.0 (brief)
\providecommand{\IndHWGal}{\ph{3.4}}           % all mergers in galaxies, galaxy catalog. MEASURED in stage 1 (known_population/galaxy_catalog/complete/galaxy_hosted): mean over realizations 100, 101, 102, 103, 105 is 3.39 (range 2.76 to 4.85; realization A 3.21). Was 3.5 (brief)
\providecommand{\IndHWAGN}{\ph{0.8}}           % all mergers in AGN, AGN catalog
\providecommand{\IndHWMix}{\ph{2.0}}           % f = 0.30, labels unknown (one population, both catalogs). MEASURED in stage 1 (known_population/both_catalogs/complete/equal_populations): mean over realizations 100, 101, 102, 103, 105 is 1.98 (over all ten 2.02, range 1.59 to 2.68; realization A 1.96). Was 1.7 (brief)
\providecommand{\IndHWMark}{\ph{1.5}}          % f = 0.30, a population for each tracer, (5 Msun, 0.10)
\providecommand{\IndHWMarkTwice}{\ph{1.4}}     % f = 0.30, a population for each tracer, (10 Msun, 0.20)
\providecommand{\IndHWKnown}{\ph{1.3}}         % f = 0.30, hosts known
\providecommand{\IndHWKnownAGNPart}{\ph{1.5}}  % f = 0.30, hosts known, the mergers in AGN alone
\providecommand{\IndHWMixTenth}{\ph{2.4}}      % f = 0.10, labels unknown
\providecommand{\IndHWKnownTenth}{\ph{2.0}}    % f = 0.10, hosts known
\providecommand{\IndCostMix}{\ph{2.4}}         % (IndHWMix/IndHWKnown)^2: mergers needed, labels unknown over hosts known. Recomputed from the stage-1 IndHWMix above: (2.0/1.3)^2 = 2.37. Was 1.7
\providecommand{\IndCostMark}{\ph{1.3}}        % (IndHWMark/IndHWKnown)^2
\providecommand{\IndCostMarkTwice}{\ph{1.2}}   % (IndHWMarkTwice/IndHWKnown)^2
\providecommand{\IndfAGNMix}{\ph{0.28^{+0.08}_{-0.08}}}  % f_AGN, f = 0.30, labels unknown, one realization

%% ---- Fig. 2(b) (figures/fig1_budget): observed-like mass spectrum (red) ----
\providecommand{\IndObsHWSpec}{\ph{10}}        % mass spectrum alone
\providecommand{\IndObsHWMix}{\ph{2.1}}        % f = 0.30, labels unknown. Estimate recomputed from the stage-1 values above (the catalog part of IndHWMix, 1/2.0^2 - 1/6.0^2, with a mass-scale part of 1/10^2: 2.08). Was 1.9
\providecommand{\IndObsHWAGN}{\ph{0.8}}        % all mergers in AGN, AGN catalog

%% ---- Fig. 3 (figures/fig2_labels): probability of an AGN host, one realization, f = 0.30, (5 Msun, 0.10) (red) ----
\providecommand{\IndProbAGNPos}{\ph{0.36}}     % mean over mergers truly in AGN, positions alone
\providecommand{\IndProbGalPos}{\ph{0.27}}     % mean over mergers truly in galaxies, positions alone
\providecommand{\IndProbAGNMark}{\ph{0.41}}    % mean over mergers truly in AGN, positions with masses and spins
\providecommand{\IndProbGalMark}{\ph{0.25}}    % mean over mergers truly in galaxies, positions with masses and spins

%% ---- Fig. 4 (figures/fig3_bias): populations inferred, f = 0.30 (red) ----
\providecommand{\IndBias}{\ph{3}}              % one population fitted, (5 Msun, 0.10): combined H0 median minus input
\providecommand{\IndBiasTwice}{\ph{6}}         % the same for (10 Msun, 0.20)
\providecommand{\IndHWOnePop}{\ph{2.6}}        % one population fitted, half-width from one catalog
\providecommand{\IndOnePopCL}{\ph{99.9}}       % credible level (percent) beyond which the input H0 lies, one population, combined
\providecommand{\IndTwoPopHzero}{\ph{67.4^{+2.1}_{-2.1}}}   % a population for each tracer, one catalog: H0
\providecommand{\IndTwoPopHW}{\ph{2.1}}                     % its half-width (open symbol of Fig. 2)
\providecommand{\IndTwoPopfAGN}{\ph{0.30^{+0.06}_{-0.06}}}  % f_AGN
\providecommand{\IndTwoPopDm}{\ph{5.1^{+1.0}_{-1.0}}}       % mass-peak difference [Msun]
\providecommand{\IndTwoPopDchi}{\ph{0.10^{+0.03}_{-0.03}}}  % mean effective spin difference
\providecommand{\IndTwoPopCL}{\ph{30}}         % credible level (percent) of the input H0, a population for each tracer, combined
\providecommand{\IndLB}{\ph{9}}                % log10 Bayes factor, a population for each tracer over one population, one catalog

%% ---- Appendix C (red) ----
\providecommand{\IndRecHzeroComb}{\ph{67.9^{+0.6}_{-0.6}}}  % equal populations, f = 0.30, labels unknown, combined: H0
\providecommand{\IndRecHzeroCL}{\ph{40}}                    % credible level (percent) of the input H0 in it
\providecommand{\IndRecfAGNComb}{\ph{0.30^{+0.03}_{-0.03}}} % the same combination: f_AGN
\providecommand{\IndRecfAGNCL}{\ph{35}}                     % credible level (percent) of the input f_AGN in it
\providecommand{\IndZerofAGNUL}{\ph{0.05}}     % no merger in AGN: 90% upper limit on f_AGN, combined

%% ---- added for the letter of BRIEF_v2 (17 Nov 2026); meanings in notes/MACROS_NEEDED.md, Sec. 5 ----
%% results (red): Sec. IV A, Table I
\providecommand{\IndAnchorHzero}{\ph{69.2^{+1.6}_{-1.6}}}   % equal populations, f = 0.30, labels unknown, population at its input, realization A: H0
\providecommand{\IndOnePopHzero}{\ph{70.9^{+2.6}_{-2.6}}}   % one population inferred from mergers whose populations differ by (5 Msun, 0.10), realization A: H0
\providecommand{\IndOnePopfAGN}{\ph{0.24^{+0.09}_{-0.09}}}  % f_AGN of the same analysis
\providecommand{\IndOnePopHzeroComb}{\ph{70.7^{+0.8}_{-0.8}}}  % one population inferred, combined over \NrealHDSingle realizations: H0 (dashed curve of Fig. 4a)
\providecommand{\IndTwoPopHzeroComb}{\ph{67.5^{+0.9}_{-0.9}}}  % a population for each tracer, combined over \NrealHD realizations: H0 (solid curve of Fig. 4a)
%% literature values (black; source: ../paper_prd/notes/PRIOR_WORK_VERIFIED.md)
\providecommand{\IndLitFifthHzero}{71.7^{+17.7}_{-11.4}}    % fifth catalog, dark and bright sirens, 90% (arXiv:2605.27227v2, Table 3)
\providecommand{\IndLitFifthNsrc}{236}                      % sources of that measurement (235 dark sirens and GW170817)
\providecommand{\IndLitFifthCatGain}{3.5}                   % percent improvement of the H0 constraint from the galaxy catalog over the mass spectrum alone
\providecommand{\IndLitQuaiafAGNUL}{0.21}                   % Veronesi et al. (Quaia): upper limit on the AGN fraction, 95%
\providecommand{\IndLitQuaiaLogLbol}{44.5}                  % log10 of the bolometric luminosity [erg/s] above which that limit holds
\providecommand{\IndLitZhuChenfAGN}{0.39^{+0.41}_{-0.32}}   % Zhu and Chen: fraction of mergers from AGN, 90%
\providecommand{\IndLitKumarHzero}{70.5^{+3.4}_{-2.9}}      % Kumar and Torres-Orjuela: H0, 68.3%, statistical
\providecommand{\IndLitKumarN}{13}                          % their number of merger-AGN associations


%% Float placement: let two full-width floats share the top of a page.
\setcounter{topnumber}{2}\setcounter{dbltopnumber}{2}\setcounter{totalnumber}{4}
\renewcommand{\topfraction}{0.9}\renewcommand{\dbltopfraction}{0.9}
\renewcommand{\textfraction}{0.08}
\renewcommand{\floatpagefraction}{0.8}\renewcommand{\dblfloatpagefraction}{0.8}

\newcommand{\kmsMpc}{\,{\rm km\,s^{-1}\,Mpc^{-1}}}
\newcommand{\Msun}{\,M_\odot}
\newcommand{\fAGN}{f_{\rm AGN}}
\newcommand{\chieff}{\chi_{\rm eff}}
\newcommand{\lB}{\log_{10}\mathcal{B}}
%% A figure file that exists is included; one that is still to be drawn leaves an empty frame.
\newcommand{\figfile}[3]{\IfFileExists{#1}{\includegraphics[width=#2]{#1}}%
  {\fbox{\parbox[c][#3][c]{\dimexpr#2-2\fboxsep-2\fboxrule\relax}{\mbox{}}}}}

\begin{document}

\title{Multi-tracer dark sirens: the Hubble constant from galaxies and active galactic nuclei with host-dependent binary black hole masses and spins}

\author{Ignacio Maga\~na Hernandez}
\affiliation{McWilliams Center for Cosmology, Department of Physics, Carnegie Mellon University, Pittsburgh, Pennsylvania 15213, USA}

\date{\today}

\begin{abstract}
We measure the Hubble constant $H_0$ from simulated binary black hole mergers hosted by two tracers of large-scale structure, galaxies and active galactic nuclei (AGN), when their masses and spins depend on the host.
Each tracer enters a multi-tracer dark siren likelihood with its own catalog and mass and spin distributions, and $H_0$ is inferred with the AGN-hosted fraction $\fAGN$.
For \MockNev{} mergers with $\fAGN=\MockfAGN$, of which the AGN-hosted have a mass peak $\MockDmuG\Msun$ higher and a mean effective spin \MockDmuChi{} higher, we find $H_0=\IndTwoPopHzero\kmsMpc$ (90\% credible interval; input \MockHzero), $\fAGN=\IndTwoPopfAGN$, and differences of $\IndTwoPopDm\Msun$ and $\IndTwoPopDchi$. % EXPECTED
Fitting one population to both tracers biases $H_0$ by $\IndBias\kmsMpc$ (90\% half-width \IndHWOnePop) and is disfavored with $\lB=\IndLB$. % EXPECTED
With known populations the half-width falls from \IndHWGal{} to \IndHWAGN{} as $\fAGN$ rises from \PriorfLo{} to \PriorfHi, because a localization volume holds \PureGalHostsPerVol{} galaxies but \PureAGNHostsPerVol{} AGN. % EXPECTED
Posteriors combined over realizations contain the inputs within 90\% credible intervals. % EXPECTED
Multi-tracer dark sirens measure $H_0$ and the population of each tracer, given an AGN catalog complete to the distances of the mergers. % EXPECTED
\end{abstract}

\maketitle

%% =====================================================================
\section{Introduction}
\label{sec:intro}

Gravitational waves from the merger of a compact binary give the luminosity distance of the source, and with a redshift they measure the Hubble constant $H_0$~\cite{Schutz:1986gp,Holz:2005df}.
Binary black hole (BBH) mergers have no generally accepted electromagnetic counterpart~\cite{Graham:2020gwr,Graham:2022xxu,Palmese:2021wcv,Veronesi:2024pjq,Cabrera:2025tzc}, so their redshifts are inferred statistically.
The dark siren method, also called the galaxy catalog method, takes the galaxies of a catalog that lie inside the localization volume of a merger as its candidate hosts~\cite{Schutz:1986gp,DelPozzo:2011vcw,LIGOScientific:2018gmd,DES:2019ccw,Gray:2019ksv,DES:2020nay,Finke:2021aom,Palmese:2021mjm}.
The spectral siren method uses the mass spectrum of the mergers, whose features are observed redshifted~\cite{Taylor:2011fs,Farr:2019twy,Mastrogiovanni:2021wsd,Ezquiaga:2022zkx}.
Current analyses infer $H_0$ and the mass spectrum together with a galaxy catalog~\cite{Mastrogiovanni:2023emh,Gray:2023wgj,Borghi:2023opd}.
With \IndLitFifthNsrc{} sources of the fifth gravitational-wave transient catalog, one of them with an electromagnetic counterpart, the LIGO--Virgo--KAGRA Collaboration finds $H_0=\IndLitFifthHzero\kmsMpc$ (90\% credible interval)~\cite{LIGOScientific:2026uyd}.
The main driver of that measurement remains the mass spectrum: the galaxy catalog improves the constraint by \IndLitFifthCatGain\% over the mass spectrum alone~\cite{LIGOScientific:2026uyd}.

A galaxy catalog adds less the more galaxies the localization volume of a typical merger holds.
A rarer class of host would confine the redshift far better.
Active galactic nuclei (AGN) are such a class, since they are a small fraction of all galaxies and the gas disk around their central black hole can bring stellar-mass black holes to merge~\cite{Bartos:2016dgn,Stone:2016wzz,Mckernan:2017ssq,Tagawa:2019osr}.
Whether a sizable fraction of the observed mergers occurs in AGN is not settled.
Searches for a spatial correlation between the mergers and AGN catalogs, at a fixed cosmology, give an upper limit of \IndLitQuaiafAGNUL{} (95\%) on the fraction hosted by AGN with bolometric luminosities above $10^{\IndLitQuaiaLogLbol}\,{\rm erg\,s^{-1}}$~\cite{Veronesi:2024tef} and, in another analysis, a fraction of $\IndLitZhuChenfAGN$ (90\%)~\cite{Zhu:2025rdv}.
AGN have been proposed as hosts for the measurement of $H_0$~\cite{Alves:2020wnl,Bom:2023zgw,Li:2025eig}, and one analysis reports $H_0=\IndLitKumarHzero\kmsMpc$ (68\%, statistical) from \IndLitKumarN{} proposed associations, with the redshift of each AGN taken as that of the merger~\cite{Kumar:2026bqt}.

Black holes that merge in AGN disks are expected to be heavier~\cite{Yang:2019okq,Yang:2019cbr,Tagawa:2021ofj} and to spin faster~\cite{Yang:2019cbr,Tagawa:2020dxe,McKernan:2021nwk,Bartos:2026xlt} than those that merge elsewhere, and analyses of the masses and spins alone have searched for such a component~\cite{Gayathri:2021xwb,Gayathri:2023bha,Li:2025iux}.
This matters for cosmology because the dark siren measurement depends on the assumed mass spectrum~\cite{Mastrogiovanni:2021wsd,Pierra:2023deu}: the mass spectrum and the catalog both constrain the redshift of each merger.
If the mergers hosted by the rare class are also the heavy ones, an analysis with one mass spectrum for all mergers assigns the wrong source-frame masses to the mergers whose redshifts the catalog constrains best.
The hosts and the mass and spin distributions then have to be inferred together.

Parts of such an analysis exist.
Bom and Palmese~\cite{Bom:2023zgw} infer $H_0$ with the fraction of mergers in AGN from an AGN catalog, with the other mergers in an unclustered background and without selection effects; Cross-Parkin \textit{et al.}~\cite{Cross-Parkin:2026pua} infer $H_0$ with a mixing fraction between two weightings of the galaxies of one catalog; and Pouw \textit{et al.}~\cite{Pouw:2026lar} infer the fraction of mergers in AGN at a fixed cosmology, with the other mergers uniform in comoving volume (see Ref.~\cite{Moncrieff:2025xle} for a forecast of the angular cross-correlation of mergers with a quasar and a galaxy catalog).
None of these analyses infers a mass or spin distribution that depends on the class of host, and all three leave that step to future work.
Mastrogiovanni \textit{et al.}~\cite{Mastrogiovanni:2023emh} write the merger rate per galaxy with a population that may depend on the properties of the galaxy and do not use that freedom.
What is missing is a dark siren analysis in which two catalogs of hosts enter on the same footing, each with its own number density, completeness and population.

In this work we present such a multi-tracer dark siren analysis.
Galaxies and AGN are two classes of host that trace the same large-scale structure, and we refer to them as tracers.
Each merger is hosted by one of the two, and each tracer enters the likelihood with its own catalog, number density, completeness, and mass and spin distributions (its population), together with the selection effects of the gravitational-wave detectors.
We infer $H_0$ with the fraction of mergers hosted by AGN and with the differences between the two populations.
Using simulated catalogs and mergers, we ask what each tracer contributes to $H_0$, what an analysis with a single population returns when the populations differ, and whether the differences can be measured with $H_0$.
The likelihood is general: any two classes of host whose mergers differ, such as galaxies of different type or stellar mass, fit it.

In Sec.~\ref{sec:method} we present the likelihood, in Sec.~\ref{sec:data} the simulated catalogs and mergers, in Sec.~\ref{sec:results} our results and in Sec.~\ref{sec:validation} their validation; we conclude in Sec.~\ref{sec:conclusions}.
We assume a flat $\Lambda$CDM cosmology with the matter density fixed at $\Omega_{\rm m}=\MockOmegam$~\cite{Planck:2015fie}, and uniform priors on $H_0$ between \PriorHzeroLo{} and $\PriorHzeroHi\kmsMpc$ and on $\fAGN$ between \PriorfLo{} and \PriorfHi.
Uncertainties are medians with 90\% equal-tailed credible intervals unless stated otherwise, and a half-width is half the width of that interval.
For a posterior combined over simulated realizations we quote the credible level of the input, the probability enclosed by the narrowest equal-tailed interval that contains it: a level below 90\% places the input inside the 90\% credible interval.

%% =====================================================================
\section{Multi-tracer likelihood}
\label{sec:method}

We use hierarchical Bayesian inference~\cite{Loredo:2004nn,Mandel:2018mve,Thrane:2018qnx,Vitale:2020aaz} to infer $H_0$, $\fAGN$ and the parameters of the populations from a set of detected mergers.
A merger is described by its source-frame primary mass $m_1$, mass ratio $q$ and effective spin $\chieff$, collected in $\theta$, by its redshift $z$ and by its sky position $\hat\Omega$, and it is hosted by one of two tracers $k$, the galaxies or the AGN.
The number of mergers per unit time at the detector, per interval of $\theta$, $z$ and $\hat\Omega$, is
\begin{equation}
\frac{dN}{d\theta\,dz\,d\hat\Omega}\propto\sum_{k}f_k\,\frac{p_{\rm pop}(\theta\,|\,\Lambda_k)\,p_k(z,\hat\Omega)}{1+z},
\label{eq:rate}
\end{equation}
where $f_{\rm AGN}$ is the fraction of mergers hosted by AGN and $f_{\rm gal}=1-\fAGN$, $p_{\rm pop}(\theta\,|\,\Lambda_k)$ is the population of tracer $k$, with parameters $\Lambda_k$, and $p_k(z,\hat\Omega)$ is the distribution of the hosts of tracer $k$ in redshift and on the sky.
The factor $1/(1+z)$ converts a merger rate per host that is constant in the source frame to the detector frame.
The host distribution is that of the galaxy catalog method~\cite{Gray:2019ksv,Gair:2022zsa,Gray:2023wgj,Mastrogiovanni:2023emh}: the sum of the cataloged objects, each with its redshift uncertainty, and of an out-of-catalog term for the hosts that the catalog misses, which follow the comoving volume with number density $n_k$ and completeness $C_k(z)$ (Appendix~\ref{app:sim}).

Given the data $d_i$ of $N$ detected mergers, the likelihood is~\cite{Mandel:2018mve}
\begin{equation}
\mathcal{L}(\{d_i\}\,|\,H_0,\fAGN,\Lambda)\propto\prod_{i=1}^{N}\frac{\sum_k f_k\,\mathcal{L}_{i,k}}{\sum_k f_k\,\beta_k}.
\label{eq:like}
\end{equation}
Here $\mathcal{L}_{i,k}$ is the likelihood of the data of merger $i$ integrated over the term of tracer $k$ in Eq.~\eqref{eq:rate}, and $\beta_k$ is the same integral with the probability of detection in place of that likelihood, the fraction of the mergers of tracer $k$ that are detected.
$H_0$ enters through the distance--redshift relation, which connects the measured distance and redshifted masses to $z$ and the source-frame masses.
Because the $\beta_k$ differ when the populations differ, $\fAGN$ is the fraction of all mergers and not of the detected ones.

The familiar analyses are limits of Eq.~\eqref{eq:like}.
Setting $\fAGN=0$ recovers the galaxy catalog method with one catalog, and $\fAGN=1$ gives the AGN catalog alone.
Replacing both $p_k$ by a distribution uniform in comoving volume recovers the mass spectrum alone, the spectral siren method; replacing $p_{\rm gal}$ alone gives an AGN catalog with an unclustered background~\cite{Bom:2023zgw,Pouw:2026lar}.
Setting $\Lambda_{\rm AGN}=\Lambda_{\rm gal}$ gives two tracers with one population, which only the positions of the mergers can tell apart.
The tracer that hosts a merger is its label, which the data do not give.
If the labels were known, Eq.~\eqref{eq:like} would split into one factor for each tracer; we use this hosts-known analysis as the reference for what the unknown labels cost.

The probability that merger $i$ is hosted by an AGN follows from Eq.~\eqref{eq:like},
\begin{equation}
r_i=\frac{\fAGN\,\mathcal{L}_{i,{\rm AGN}}}{\fAGN\,\mathcal{L}_{i,{\rm AGN}}+(1-\fAGN)\,\mathcal{L}_{i,{\rm gal}}},
\label{eq:member}
\end{equation}
which we evaluate at the posterior median of the parameters.
The integrals are Monte Carlo sums over posterior samples of each merger and over a set of simulated signals (injections), and we require the effective number of injections to exceed \NeffFactor{} per detected merger~\cite{Farr:2019rap} (Appendix~\ref{app:mc}).

%% =====================================================================
\section{Simulated catalogs and mergers}
\label{sec:data}

\begin{figure*}
\figfile{figures/fig_concept.pdf}{\textwidth}{0.4\textwidth}
\caption{The simulated catalogs and mergers.
(a) A slice of the galaxy catalog (blue points; \IndConceptNGal{} objects) and of the AGN catalog (orange diamonds; \IndConceptNAGN{} objects) of one realization, within $\IndConceptDec^\circ$ of the equator and out to $z=\IndConceptzMax$, drawn in comoving distance; the bar is \IndConceptScale{} Mpc long.
The black outline is the 90\% localization region of one AGN-hosted merger at $z=\IndConceptzTrue$, enlarged in the inset; it holds \IndConceptVolNGal{} galaxies and \IndConceptVolNAGN{} AGN, and the circle marks the true host of the merger.
(b) Distribution of the source-frame primary mass $m_1$ of the galaxy-hosted (blue; \IndConceptNgalHosted{} mergers) and AGN-hosted (orange; \IndConceptNAGNHosted{} mergers) mergers detected in a realization whose populations depend on the host; the dashed lines mark the input locations of the two mass peaks, \MockMuG{} and $\MockMuGAGN\Msun$.
(c) The same for the effective spin $\chieff$, with dashed lines at the two input means, \MockMuChi{} and \MockDmuChi.
The localization region holds far fewer AGN than galaxies, and the AGN-hosted mergers are heavier and spin faster.}
\label{fig:concept}
\end{figure*}

\begin{figure}
\figfile{figures/fig_anchor.pdf}{\columnwidth}{0.8\columnwidth}
\caption{Joint posterior of $H_0$ and the fraction of mergers hosted by AGN, $\fAGN$, for mergers with one population common to both tracers and held at its input.
The dashed line is the 90\% credible region of one realization of \MockNev{} mergers, the filled region with a solid outline is that of the combination of \NrealKnown{} realizations, and the black plus sign is the input.
Both parameters are measured together, and the 90\% credible region of the combination contains the input.}
\label{fig:anchor}
\end{figure}

The likelihood of Eq.~\eqref{eq:like} makes three demands on simulated data, and we build the simulation to meet each of them.
First, the two tracers must trace one density field.
We draw a lognormal matter field~\cite{Tessore:2023zyk} in shells \MockShellWidth{} Mpc thick out to $z=\MockzCat$, with the matter power spectrum~\cite{Lewis:1999bs} of a flat $\Lambda$CDM cosmology with $H_0=\MockHzero\kmsMpc$ and $\Omega_{\rm m}=\MockOmegam$~\cite{Planck:2015fie}, and populate it with two catalogs.
The galaxies have a comoving number density of $\MocknGal\,{\rm Mpc^{-3}}$, that of galaxies brighter than \MockLcut{} times the characteristic luminosity of the $B$-band luminosity function adopted for the GLADE catalogs~\cite{Dalya:2018cnd,Dalya:2021ewn}, and a linear bias of \MockbGal.
The AGN have a number density of $\MocknAGN\,{\rm Mpc^{-3}}$, representative of AGN with 2--10 keV luminosities above $10^{\MockLXcut}\,{\rm erg\,s^{-1}}$, a class selected in hard X-ray surveys~\cite{Powell:2018}, and a bias of \MockbAGN.
The catalogs hold $\MockNGal$ galaxies and $\MockNAGN$ AGN (Fig.~\ref{fig:concept}a).
Each object has a Gaussian redshift error of $\MockDz(1+z)$, which is realized in the cataloged redshift, so that the redshift uncertainty in the likelihood is the true one.
Both catalogs are complete to $z=\MockzCat$, beyond the most distant detected merger at $z=\MockzHorizon$; the out-of-catalog term then vanishes and no number density enters the likelihood.

\begin{figure*}
\figfile{figures/fig1_budget.pdf}{\textwidth}{0.4\textwidth}
\caption{The 90\% half-width of $H_0$ from \MockNev{} mergers against the fraction of mergers hosted by AGN, $\fAGN$, with the populations held at their inputs: (a) for the fiducial spectrum, with \MockPeakFrac{} of the mergers in the peak, and (b) for a spectrum like the observed one, with \ObsPeakFrac{} in the peak.
Circles and solid lines: labels unknown, with one population for both tracers.
Triangles and dotted lines: hosts known.
Squares and diamonds with dashed lines: a population for each tracer, for mergers whose populations differ by $\Delta m=\MockDmuG\Msun$ and $\Delta\chi=\MockDmuChi$ and by twice those values.
Open square: a population for each tracer inferred with $H_0$.
Horizontal lines: the mass spectrum alone (long dashes) and the galaxy catalog alone at $\fAGN=0$ (short dashes).
Points are means over \IndNreal{} realizations, the bars span the full range over those realizations, and the points are displaced horizontally for clarity.
The half-width falls steeply with $\fAGN$, and masses and spins that depend on the host move it toward the value for hosts known.} % EXPECTED
\label{fig:budget}
\end{figure*}

Second, the mergers must follow Eq.~\eqref{eq:rate}.
A merger is hosted by an AGN with probability $\fAGN$ and by a galaxy otherwise, every object of a tracer is equally likely to host it, and the merger rate per host is constant in the source frame.
We simulate $\fAGN=\IndfList$; at $\fAGN=\MockfAGN$ an AGN hosts mergers at \MockRateRatio{} times the rate of a galaxy.
The primary mass follows a power law with a Gaussian peak~\cite{Talbot:2018cva}.
Our fiducial spectrum has \MockPeakFrac{} of the mergers in a peak at $\MockMuG\Msun$ of width $\MockSigmaG\Msun$ and the rest in a power law of index $-\MockAlpha$ between \MockMmin{} and $\MockMmax\Msun$, a sharp mass scale chosen to make the mass spectrum as informative about redshift as it can be.
A second spectrum is like the observed one~\cite{KAGRA:2021duu}, with \ObsPeakFrac{} of the mergers in a peak at $\ObsMuG\Msun$ and a power law of index $-\ObsAlpha$; \ObsDetPeakShare{} of its detected mergers are in the peak, against \MockDetPeakShare{} for the fiducial spectrum.
The effective spin is Gaussian with mean \MockMuChi{} and width \MockSigmaChi.
When the populations depend on the host, the AGN-hosted mergers have the mass peak higher by $\Delta m=\MockDmuG\Msun$ and the mean effective spin higher by $\Delta\chi=\MockDmuChi$ (Fig.~\ref{fig:concept}b and c); we also simulate differences twice as large, $\IndDmTwice\Msun$ and \IndDchiTwice.

Third, detection must be a function of the data alone, because Eq.~\eqref{eq:like} is a probability density of the observed data of detected mergers~\cite{Essick:2023upv}.
A merger is detected when its observed signal-to-noise ratio exceeds \MockSNRthr, and the uncertainties of its masses, spin, distance and sky position scale inversely with that observed value~\cite{Farah:2023vsc} (Appendix~\ref{app:sim}).
The median fractional uncertainty of the luminosity distance is \MockSigLnDL\% and the median redshift of the detected mergers is \MockzMedian.
The 90\% localization volume of a galaxy-hosted merger then holds a median of \PureGalHostsPerVol{} galaxies, and that of an AGN-hosted merger a median of \PureAGNHostsPerVol{} AGN.
A realization consists of one density field with its two catalogs, \MockNev{} detected mergers with \ProdNsamp{} posterior samples each (\MCNsampHi{} when every merger is hosted by an AGN), and its own injection set for the selection effects.

%% =====================================================================
\section{Results}
\label{sec:results}

\begin{table*}
\caption{Results by analysis for one realization of \MockNev{} mergers with $\fAGN=\MockfAGN$ and $H_0=\MockHzero\kmsMpc$; where the populations depend on the host the inputs are $\Delta m=\MockDmuG\Msun$ and $\Delta\chi=\MockDmuChi$.
Entries are medians with 90\% credible intervals.
$\lB$ is the Bayes factor against the analysis with one population inferred from the same mergers, with uniform priors on $\Delta m$ between $\PriorDmuGLo$ and $\PriorDmuGHi\Msun$ and on $\Delta\chi$ between $\PriorDmuChiLo$ and $\PriorDmuChiHi$.}
\label{tab:results}
\begin{ruledtabular}
\begin{tabular}{lccccc}
Analysis & $H_0$ [${\rm km\,s^{-1}\,Mpc^{-1}}$] & $\fAGN$ & $\Delta m$ [$M_\odot$] & $\Delta\chi$ & $\lB$\\
\hline
\multicolumn{6}{l}{\textit{Mergers with a population common to both tracers}}\\
Population held at its input & $\IndAnchorHzero$ & $\IndfAGNMix$ & 0 (fixed) & 0 (fixed) & \\
\hline
\multicolumn{6}{l}{\textit{Mergers whose populations depend on the host}}\\
One population inferred & $\IndOnePopHzero$ & $\IndOnePopfAGN$ & 0 (fixed) & 0 (fixed) & 0\\
A population for each tracer inferred & $\IndTwoPopHzero$ & $\IndTwoPopfAGN$ & $\IndTwoPopDm$ & $\IndTwoPopDchi$ & $\IndLB$\\
\end{tabular}
\end{ruledtabular}
\end{table*}

\begin{figure*}
\figfile{figures/fig_populations.pdf}{\textwidth}{0.4\textwidth}
\caption{Mergers whose populations depend on the host, with $\Delta m=\MockDmuG\Msun$, $\Delta\chi=\MockDmuChi$ and $\fAGN=\MockfAGN$.
(a) Posterior density of $H_0$ combined over realizations, with one population fitted to all mergers (dashed; \NrealHDSingle{} realizations) and with a population for each tracer (solid; \NrealHD{} realizations); the shaded bands are the 90\% credible intervals and the vertical dashed line is the input.
(b) The 90\% credible region of the difference of the mass-peak location, $\Delta m$, and of the difference of the mean spin, $\Delta\chi$, from one realization with a population for each tracer; the cross is the input and the circle marks equal populations.
(c) Distribution over the mergers of one realization of the probability of an AGN host, Eq.~\eqref{eq:member}, for mergers truly in galaxies (filled) and truly in AGN (dark line), from positions alone (top) and from positions with masses and spins (bottom); the triangles mark the means and the dashed line marks $\fAGN$.
One population fitted to mergers of two kinds displaces $H_0$ by $\IndBias\kmsMpc$, and a population for each tracer removes the displacement.} % EXPECTED
\label{fig:populations}
\end{figure*}

\subsection{A population common to both tracers}
\label{sec:anchor}

We first analyze mergers whose masses and spins do not depend on their host, with $\fAGN=\MockfAGN$, and hold the population at its input, so that $H_0$ and $\fAGN$ are the only free parameters.
This analysis is given the true population, and it tests the two host distributions and the selection effects of Eq.~\eqref{eq:like} alone.
From one realization of \MockNev{} mergers we find $H_0=\IndAnchorHzero\kmsMpc$ and $\fAGN=\IndfAGNMix$ (Fig.~\ref{fig:anchor} and Table~\ref{tab:results}), and both inputs lie inside the 90\% credible intervals.
The combination of \NrealKnown{} realizations gives $H_0=\IndRecHzeroComb\kmsMpc$ and $\fAGN=\IndRecfAGNComb$, again with both inputs inside the 90\% credible intervals (Sec.~\ref{sec:validation}).
The mergers thus tell apart two catalogs that trace the same structure, and the fraction hosted by each is measured with $H_0$.
The limit of this analysis with no merger in an AGN is the familiar one with a galaxy catalog, to which we turn next.

\subsection{The information in each tracer}
\label{sec:budget}

The two tracers contribute very differently to $H_0$.
In this subsection the populations are held at their inputs and the half-widths are means over \IndNreal{} realizations of \MockNev{} mergers.
When no merger is hosted by an AGN, the galaxy catalog gives a 90\% half-width of $\IndHWGal\kmsMpc$, to be compared with \IndHWSpec{} from the mass spectrum of the same mergers with no catalog. % EXPECTED
A catalog with \PureGalHostsPerVol{} candidate hosts in a localization volume thus roughly halves the half-width that the mass scale gives, against the improvement of \IndLitFifthCatGain\% found for the fifth catalog~\cite{LIGOScientific:2026uyd}; our galaxy catalog is complete to the distances of the mergers. % EXPECTED
When every merger is hosted by an AGN, the AGN catalog gives $\IndHWAGN\kmsMpc$, with \PureAGNHostsPerVol{} candidate hosts in a localization volume.

Figure~\ref{fig:budget}(a) shows the half-width for $\fAGN=\IndfList$.
With the labels unknown it is $\IndHWMixTenth\kmsMpc$ at $\fAGN=\MockfAGNTenth$ and \IndHWMix{} at $\fAGN=\MockfAGN$. % EXPECTED
A minority of AGN-hosted mergers carries most of the constraint. % EXPECTED

Part of this information is lost because the labels are unknown.
With the hosts known the half-widths would be \IndHWKnownTenth{} and $\IndHWKnown\kmsMpc$, and at $\fAGN=\MockfAGN$ the AGN-hosted mergers alone would give \IndHWKnownAGNPart. % EXPECTED
For a half-width that scales as the inverse square root of the number of mergers, the analysis with labels unknown needs \IndCostMix{} times as many mergers as the analysis with hosts known. % EXPECTED
The likelihood finds the AGN-hosted mergers statistically, through Eq.~\eqref{eq:member}, and we argue that positions identify them weakly because galaxies and AGN trace the same structure, so that a localization volume that holds AGN also holds galaxies. % EXPECTED

\subsection{Populations that depend on the host}
\label{sec:populations}

We now let the AGN-hosted mergers be heavier and spin faster, with $\Delta m=\MockDmuG\Msun$, $\Delta\chi=\MockDmuChi$ and $\fAGN=\MockfAGN$, and infer the populations with $H_0$.
Heavier mergers are detected to larger distances, so \IndConceptNAGNHosted{} of the \MockNev{} detected mergers of one realization are AGN-hosted.

First, we fit one population to both tracers, as current dark siren analyses do: the location and width of the mass peak and the mean and width of the effective spin distribution are inferred with $H_0$ and $\fAGN$.
In the posterior combined over \NrealHDSingle{} realizations the median of $H_0$ lies $\IndBias\kmsMpc$ above the input (Fig.~\ref{fig:populations}a), to be compared with a 90\% half-width of \IndHWOnePop{} from one realization of \MockNev{} mergers, and the input is excluded at the \IndOnePopCL\% credible level. % EXPECTED
For differences twice as large the median lies $\IndBiasTwice\kmsMpc$ above the input. % EXPECTED
We read the bias as an error in the mass scale: the single mass peak lies below that of the AGN-hosted mergers, which are then assigned redshifts that are too high for their distances, and these are the mergers whose redshifts the AGN catalog constrains best. % EXPECTED
We have not separated the effect of the mass difference from that of the spin difference.

Second, we give each tracer its own population: the mass peak of the AGN-hosted mergers is displaced by $\Delta m$ and their mean effective spin by $\Delta\chi$, and both differences are inferred.
From one realization we find $H_0=\IndTwoPopHzero\kmsMpc$, $\fAGN=\IndTwoPopfAGN$, $\Delta m=\IndTwoPopDm\Msun$ and $\Delta\chi=\IndTwoPopDchi$ (Table~\ref{tab:results}), with each input inside its 90\% credible interval, and equal populations lie outside the 90\% credible region of the two differences (Fig.~\ref{fig:populations}b). % EXPECTED
The Bayes factor of this analysis against the one with a single population is $\lB=\IndLB$, where $\lB>0$ favors a population for each tracer and a value of this size is decisive. % EXPECTED
The half-width of $H_0$ is $\IndTwoPopHW\kmsMpc$ (open square in Fig.~\ref{fig:budget}a), against \IndHWOnePop{} with one population: a population for each tracer removes the bias and is also more precise. % EXPECTED

Finally, the masses and spins help to identify the hosts.
For mergers truly in AGN, the mean probability of an AGN host is \IndProbAGNPos{} from positions alone and \IndProbAGNMark{} from positions with masses and spins; for mergers truly in galaxies it is \IndProbGalPos{} and \IndProbGalMark{} (Fig.~\ref{fig:populations}c), where a merger with no information would have $\fAGN=\MockfAGN$. % EXPECTED
With the populations held at their inputs, the half-width of $H_0$ is $\IndHWMark\kmsMpc$ for these mergers and \IndHWMarkTwice{} for differences twice as large, between the values for labels unknown and for hosts known (Fig.~\ref{fig:budget}a), and the number of mergers needed relative to hosts known falls from \IndCostMix{} to \IndCostMark{} and \IndCostMarkTwice. % EXPECTED
We read this gain as the masses and spins sharpening the assignment of mergers to tracers. % EXPECTED
The comparison is not between the same mergers, however: more of the detected mergers are AGN-hosted when these are heavier, which also lowers the half-width.

%% =====================================================================
\section{Validation}
\label{sec:validation}

We validate these results with sets of realizations drawn from known inputs.
Specifically, we want to test whether the multi-tracer likelihood can spuriously assign mergers to AGN that host none, whether a recovery of the inputs in one realization could be a chance agreement, and whether the results rest on the sharp mass scale of the fiducial spectrum.

First, we combine the posteriors of independent realizations, each with its own density field, catalogs, mergers and injections.
For a population common to both tracers (Sec.~\ref{sec:anchor}), the input $H_0$ lies at the \IndRecHzeroCL\% credible level of the posterior combined over \NrealKnown{} realizations and the input $\fAGN$ at the \IndRecfAGNCL\% level.
For populations that depend on the host and a population for each tracer, the posterior combined over \NrealHD{} realizations gives $H_0=\IndTwoPopHzeroComb\kmsMpc$ with the input at the \IndTwoPopCL\% level, whereas with one population the combination of \NrealHDSingle{} realizations gives $\IndOnePopHzeroComb\kmsMpc$ and excludes the input at the \IndOnePopCL\% level (Fig.~\ref{fig:populations}a). % EXPECTED
The combined interval, $H_0=\IndRecHzeroComb\kmsMpc$ in the first case against $\IndAnchorHzero\kmsMpc$ from one realization, would reveal a bias well below the half-width of one realization.

Second, we analyze mergers that are all hosted by galaxies with both catalogs and $\fAGN$ free.
The posterior of $\fAGN$ has its maximum at zero, and from \IndNreal{} realizations combined we find $\fAGN<\IndZerofAGNUL$ at 90\% credibility. % EXPECTED
At both ends of Fig.~\ref{fig:budget}(a), where every merger is hosted by one tracer, the analysis with both catalogs gives nearly the half-width of the analysis with the single catalog. % EXPECTED
The likelihood does not find AGN-hosted mergers where there are none. % EXPECTED

Third, we replace the fiducial mass spectrum by the one like the observed spectrum (Fig.~\ref{fig:budget}b).
The mass spectrum alone then gives a half-width of $\IndObsHWSpec\kmsMpc$, against \IndHWSpec{} for the fiducial spectrum, while the analysis with labels unknown at $\fAGN=\MockfAGN$ gives \IndObsHWMix{} against \IndHWMix{} and the AGN catalog with every merger in an AGN gives \IndObsHWAGN{} against \IndHWAGN. % EXPECTED
The constraint from the AGN-hosted mergers hardly uses the mass scale, so the results of Sec.~\ref{sec:budget} hold for a realistic spectrum, in which the catalogs matter more; the bias of Sec.~\ref{sec:populations} we have measured for the fiducial spectrum only. % EXPECTED

Finally, the quoted intervals carry a Monte Carlo error, because the likelihood is evaluated with sums over posterior samples and injections (Appendix~\ref{app:mc}).
With \ProdNsamp{} posterior samples per merger, \MCMixedNEightK{} redraws of the samples of one realization with $\fAGN=\MockfAGN$ move the median of $H_0$ by $\MCMixedHzeroSdEightK\kmsMpc$ (standard deviation) and a new injection set moves it by \MCMixedInjHzeroSd; when every merger is hosted by an AGN, the redraws move it by \MCAGNHzeroSdEightK, and by \MCAGNHzeroSdThirtyTwoK{} at the \MCNsampHi{} samples per merger with which these sets are analyzed.
These shifts are to be compared with half-widths of \IndHWMix{} and $\IndHWAGN\kmsMpc$. % EXPECTED
They are independent between realizations, which have their own samples and injections, and average down in the combined posteriors.

%% =====================================================================
\section{Discussion and conclusions}
\label{sec:conclusions}

We have presented a multi-tracer dark siren analysis, in which galaxies and AGN enter one likelihood with their own catalogs and populations, and applied it to simulated realizations of \MockNev{} BBH mergers.
We find that $H_0$ and the fraction of mergers hosted by AGN are measured together, that the constraint on $H_0$ comes mostly from the AGN-hosted mergers, and that an analysis with one population is biased by $\IndBias\kmsMpc$ when the AGN-hosted mergers are heavier by $\MockDmuG\Msun$ and spin faster by \MockDmuChi, while a population for each tracer recovers $H_0$ and measures both differences. % EXPECTED

These results rest on four inputs.
First, the measurement of $H_0$ needs an AGN catalog that is complete over the volume of the detected mergers, here $z<\MockzHorizon$.
The searches for a spatial correlation find that the completeness of the AGN catalog limits their constraint, because the fraction and the completeness enter their likelihood as a product~\cite{Veronesi:2023ugk,Veronesi:2024tef,Zhu:2025rdv}, and Pouw \textit{et al.} find from simulated catalogs that the catalog must be complete within the localization volumes~\cite{Pouw:2026lar}.
Second, $\fAGN$ needs the abundance of AGN.
The mergers constrain the merger rate per AGN relative to that per galaxy, $R$, and $\fAGN/(1-\fAGN)=R\,N_{\rm AGN}/N_{\rm gal}$, with $N_k$ the number of hosts of tracer $k$.
With complete catalogs the $N_k$ are the catalog counts and no assumed density enters; from one realization we find $R=\NormRateRatioCompleteProd$, an interval that contains the input of \MockRateRatio.
With catalogs limited in flux so that they hold the hosts of \MockEventCompEighteen\% of the detected mergers, an AGN number density assumed at twice its input moves $\fAGN$ from $\NormfEighteenProd$ to $\NormfEighteenProdTwo$ and $H_0$ from $\NormHzeroEighteenProd$ to $\NormHzeroEighteenProdTwo\kmsMpc$; Pouw \textit{et al.} take this density from a published quasar luminosity function~\cite{Pouw:2026lar}.
Third, the size of the population differences is not known, and the bias of the single-population analysis grows with it. % EXPECTED
Pouw \textit{et al.} find their position-based results inconsistent with a spin-based ranking of candidate AGN mergers and conclude that the two must be inferred jointly~\cite{Pouw:2026lar}; an analysis with a population for each tracer measures the differences and does not need them in advance.
Fourth, we assumed the form of the mass spectrum and inferred the location and width of its peak, while the dark siren $H_0$ is known to depend on that form~\cite{Mastrogiovanni:2021wsd,Pierra:2023deu}.

In the near term the same likelihood applies to the mergers of the fifth catalog with all-sky catalogs of galaxies and quasars~\cite{Dalya:2021ewn,Storey-Fisher:2023gca}, and a companion paper in preparation treats catalogs limited in flux and number densities inferred with $H_0$.
The step change is a set of \MockNev{} mergers, several times the \IndLitFifthNsrc{} sources of the fifth catalog, which we expect from the coming observing runs at improved sensitivity~\cite{KAGRA:2013rdx}, together with an AGN catalog complete to their distances.
The systematics that remain open are the AGN missing from catalogs, among them the obscured ones, the abundance of the AGN that host mergers, the evolution of $\fAGN$ and of the populations with redshift, and the Monte Carlo accuracy of the likelihood for sets of this size.
The differences $\Delta m$ and $\Delta\chi$ that multi-tracer dark sirens measure with $H_0$ are quantities that models of black hole mergers in AGN disks predict, so each such measurement of the expansion rate is also a test of those models.

\appendix

%% =====================================================================
\section{Host distributions, simulated measurements and priors}
\label{app:sim}

The host distribution of tracer $k$ in Eq.~\eqref{eq:rate} is
\begin{equation}
Z_k\,p_k(z,\hat\Omega)=\sum_{j\in k}K_j(z)\,\frac{\Theta_j(\hat\Omega)}{\Omega_{\rm pix}}+n_k\,[1-C_k(z)]\,\frac{dV_c}{dz\,d\hat\Omega},
\label{eq:hosts}
\end{equation}
where the sum runs over the cataloged objects, $K_j(z)$ is a Gaussian at the cataloged redshift of object $j$ with its redshift error, multiplied by the comoving volume element and normalized to unity, $\Theta_j$ is one in the sky pixel of object $j$ and zero elsewhere, the pixels being equal-area ones of solid angle $\Omega_{\rm pix}$ with resolution parameter $N_{\rm side}=\MockNside$, and $V_c$ is the comoving volume.
The constant $Z_k$ normalizes $p_k$ to unity over the sky and over $0<z<\NormZmaxProd$, a range that contains every detected merger and injection; for a complete catalog it is the number of cataloged objects.

The optimal signal-to-noise ratio of a merger of redshifted chirp mass $\mathcal{M}_z$ at luminosity distance $d_L$ is
\begin{equation}
\rho_{\rm opt}=\MockRhoRef\,\Big(\frac{\mathcal{M}_z}{\MockMcRef\Msun}\Big)^{5/6}\,\frac{1\,{\rm Gpc}}{d_L},
\label{eq:snr}
\end{equation}
and the observed one is $\rho_{\rm obs}=\rho_{\rm opt}+n$ with $n$ a standard normal variate; a merger is detected if $\rho_{\rm obs}\geq\MockSNRthr$~\cite{Farah:2023vsc,Essick:2023upv}.
The measurement uncertainties are
\begin{equation}
\big(\sigma_{\ln\mathcal{M}_z},\,\sigma_{\ln q},\,\sigma_{\chieff}\big)=\big(\MockSigLnMc,\,\MockSigLnq,\,\MockSigChi\big)\,\frac{\MockSNRthr}{\rho_{\rm obs}},
\label{eq:widths}
\end{equation}
and the sky position has a Gaussian uncertainty between $\MockSkyMin^\circ$ and $\MockSkyMax^\circ$ for the detected mergers.
The distance is measured through $\rho_{\rm obs}$ and $\mathcal{M}_z$ in Eq.~\eqref{eq:snr}.
The posterior samples of each merger are drawn from the exact posterior of these measurements under a prior that is divided out in the likelihood.
The injections that give the $\beta_k$ are subject to the same detection rule: of $\MockNinjDrawn$ signals drawn for a realization, $\MockNinjFound$ are detected.

The priors are uniform.
Besides those on $H_0$ and $\fAGN$ (Sec.~\ref{sec:intro}), they run from \PriorMuGLo{} to $\PriorMuGHi\Msun$ for the location of the mass peak, from \PriorSigmaGLo{} to $\PriorSigmaGHi\Msun$ for its width, from $\PriorMuChiLo$ to $\PriorMuChiHi$ for the mean effective spin and from \PriorSigmaChiLo{} to \PriorSigmaChiHi{} for its width, with the ranges of $\Delta m$ and $\Delta\chi$ given in Table~\ref{tab:results}.
The other parameters of the mass spectrum are held at their inputs.

%% =====================================================================
\section{Monte Carlo evaluation of the likelihood}
\label{app:mc}

We estimate the logarithm of Eq.~\eqref{eq:like} as
\begin{equation}
\ln\hat{\mathcal{L}}=\sum_{i=1}^{N}\Big[\ln\bar w_i+\frac{\sigma_i^2}{2}\Big]-N\ln\hat\beta+\frac{N(N+3)}{2N_{\rm eff}},
\label{eq:estimator}
\end{equation}
where $\bar w_i$ is the mean over the posterior samples of merger $i$ of the ratio of the rate of Eq.~\eqref{eq:rate} to the prior of the samples, $\sigma_i^2$ is the variance of $\ln\bar w_i$, $\hat\beta$ is the corresponding mean over the detected injections, and $N_{\rm eff}$ is their effective number~\cite{Farr:2019rap}.
The term $\sigma_i^2/2$ corrects the logarithm of a Monte Carlo mean, which is low by half its variance on average~\cite{Essick:2022ojx,Talbot:2023pex}, and the last term accounts for the uncertainty of $\hat\beta$~\cite{Farr:2019rap}.
The posteriors of $H_0$ and $\fAGN$ with the populations held at their inputs are computed on grids, and the analyses that infer the populations use nested sampling~\cite{Speagle:2019ivv}.

\bibliography{refs}

\end{document}
